% Do not change document class, margins, fonts, etc.
\documentclass[a4paper,oneside,bibliography=totoc]{scrartcl}

% some useful packages (add more as needed)
\usepackage[utf8]{inputenc}
\usepackage{graphicx}
\usepackage{latexsym}
\usepackage{amsmath}
\usepackage{amssymb}
\usepackage{tabularx}
\usepackage{booktabs}
\usepackage{algorithm} % you can modify the algorithm style to your liking
\usepackage{algorithmic}
\usepackage{csquotes}
\renewcommand{\algorithmiccomment}[1]{\hfill\textit{// #1}}
\usepackage[usenames,dvipsnames]{xcolor}
\usepackage[colorlinks,citecolor=Green]{hyperref} % you may change/remove the colors
\usepackage{lipsum} % you do not need this
\usepackage{environ}
\usepackage{tikz}
\usetikzlibrary{positioning,shapes.gates.logic.US,shapes.gates.logic.IEC,automata,arrows, backgrounds}
\usepackage{natbib}
\bibliographystyle{chicagoa}
\setcitestyle{authoryear,round,semicolon,aysep={},yysep={,}} \let\cite\citep

% example enviroments and commands (add more as needed)
\newtheorem{definition}{Definition} \newtheorem{proposition}{Proposition}

\newcommand\todo[1]{\textcolor{red}{\framebox{\footnotesize TODO}\expandafter\ifx\expandafter\relax\detokenize{#1}\relax\else\ #1\fi}}


% DO NOT MODIFY
\newlength\oldparskip
\NewEnviron{answer}[1][8cm]{%
  \sbox{0}{\parbox[t][#1][t]{0.99\linewidth}{%
      \setlength{\oldparskip}{\parskip}%
      \setlength{\parskip}{.5em}%
      \BODY%
      ~%
      \setlength{\parskip}{\oldparskip}%
    }}%

  \begin{tikzpicture}
    % inner sep restores the \fbox padding
    \node[inner sep=\fboxsep, outer sep=0pt] (A) {\usebox{0}};
    % border thickness like \fboxrule
    \draw[line width=\fboxrule] (A.south west) rectangle (A.north east);
  \end{tikzpicture}%
  \vskip\baselineskip
}

\newcommand{\answerfig}[2]{%
  \par\noindent
  \makebox[\linewidth][c]{\includegraphics[width=0.98\linewidth]{#1}}\par
  \vspace{-0.4\baselineskip}
  {\footnotesize\centering #2\par}
  \vspace{0.25em}
}


\begin{document}

\title{ML26, Assignment 2}
\author{Selma Nezihoglu (snezihog, 2390917)\\
  Abhaykumar Narendrakumar Monasara (amonasar, S2188510)}
\date{\today}

\maketitle

\section*{Task 1a) and Task 1c)}
\begin{answer}[19.5cm]
  \small
  \answerfig{task1_densities.png}{KDEs of all 57 features: (a) raw data (provided plot),
  (b) raw word frequencies up to their 99th percentile, (c) z-scored training data on $[-1,5]$.}

  \textbf{1a)} The provided plot (a) is almost uninformative: all curves collapse into a single spike at $0$.
  The $x$-axis has to cover the largest value, \texttt{capital\_run\_length\_total} (up to $15841$), whereas
  word and character frequencies never exceed $42.8$ and $32.5$, respectively. On one axis spanning such
  different scales, no individual distribution is visible. Plotting the raw word frequencies only up to
  their 99th percentile (b) shows what plot (a) hides:
  \begin{itemize}
    \item \emph{Sparsity.} For the median feature, $84\%$ of the training emails have the value exactly $0$.
    \item \emph{Right skew.} The distributions are strongly right-skewed (median skewness $9.4$): almost all
          mass lies near $0$, followed by a long tail formed by a few emails with large values.
    \item \emph{Exceptions.} Only the three capital-run-length features are never $0$, and they also have by
          far the largest range.
  \end{itemize}

  \textbf{1c)} We z-scored both training and test data using the mean and standard deviation of the
  training data. All features now share one scale, so every density is visible (c), and on the training data
  each feature has mean $0$ and variance $1$. Because z-scoring is an affine map with positive slope, it
  changes only location and scale, not shape: each curve is still a sharp peak with a long right tail. The
  peak lies slightly \emph{below} $0$, since the many zero entries are mapped to $-\mu_j/\sigma_j<0$. Three
  observations stand out:
  \begin{itemize}
    \item \emph{Extreme outliers.} Z-scores range from $-0.93$ to $45.1$, although $98.4\%$ of all values lie
          in $[-1,3]$. The largest values belong to \texttt{capital\_run\_length\_longest} ($45.1$),
          \texttt{capital\_run\_length\_average} ($40.3$) and \texttt{char\_freq\_\#} ($38.1$); in each case a
          single email lies more than $38$ standard deviations above the mean.
    \item \emph{Asymmetry.} Since the raw features are non-negative, the z-scores of feature $j$ are bounded
          below by $-\mu_j/\sigma_j$, and no value falls below $-0.93$. The right tails, by contrast, are very
          long, so the features are far from Gaussian.
    \item \emph{Test data is not exactly standardized.} Test means lie in $[-0.06,0.09]$ and test variances in
          $[0.05,2.76]$. This is expected: the test data is transformed with the training statistics, so no
          information from the test set enters preprocessing.
  \end{itemize}
\end{answer}
\section*{Task 2a)}
\begin{answer}[19.5cm]
  \small
  \textbf{Claim.} Consider $p(y=1\mid\mathbf x)=\sigma(\mathbf w^\top\mathbf x+b)$ with empirical risk
  \[
    \hat R(\mathbf w,b)=\frac1N\sum_{i=1}^N \ell(y_i,s_i),\quad s_i=\mathbf w^\top\mathbf x_i+b,\quad
    \ell(y,s)=-y\log\sigma(s)-(1-y)\log\sigma(-s).
  \]
  Rescale and shift every feature, $\tilde x_{ij}=a_j x_{ij}+c_j$ with constants $a_j\neq0$, $c_j\in\mathbb R$.
  Then the smallest achievable empirical risk (the infimum, which is a minimum whenever a minimizer exists)
  is the same on $\tilde{\mathbf X}$ and on $\mathbf X$.

  \textbf{Proof.} The loss of example $i$ depends on the parameters only through its score $s_i$. Parameters
  that yield the same score for every example therefore have the same empirical risk. Given any
  $(\mathbf w,b)$ for the original data, set
  \[
    \tilde w_j=\frac{w_j}{a_j},\qquad \tilde b=b-\sum_{j=1}^D\frac{w_j c_j}{a_j}.
  \]
  Then for every example $i$
  \[
    \tilde{\mathbf w}^\top\tilde{\mathbf x}_i+\tilde b
    =\sum_{j=1}^D\frac{w_j}{a_j}(a_j x_{ij}+c_j)+b-\sum_{j=1}^D\frac{w_j c_j}{a_j}
    =\mathbf w^\top\mathbf x_i+b,
  \]
  hence $\hat R_{\tilde{\mathbf X}}(\tilde{\mathbf w},\tilde b)=\hat R_{\mathbf X}(\mathbf w,b)$. The map
  $(\mathbf w,b)\mapsto(\tilde{\mathbf w},\tilde b)$ is a bijection with inverse $w_j=a_j\tilde w_j$,
  $b=\tilde b+\sum_j\tilde w_j c_j$. Both problems therefore attain exactly the same set of risk values, and
  their infima coincide. In particular, if $(\mathbf w^*,b^*)$ is an ERM estimate on $\mathbf X$, its image
  $(\tilde{\mathbf w}^*,\tilde b^*)$ is an ERM estimate on $\tilde{\mathbf X}$ with the same risk and the
  same predictions. \hfill$\square$

  The bias term is essential for the shift: without $b$, the term $\sum_j w_j c_j/a_j$ cannot be absorbed.
  Pure rescaling ($c_j=0$) still leaves the minimal risk unchanged, but a shift in general does not.

  \textbf{Why z-scores, then?} The ERM solutions are equivalent, but \emph{finding} them is not equally easy:
  \begin{itemize}
    \item \emph{Optimization.} Gradient descent uses a single step size for all weights. On the raw data the
          features differ by orders of magnitude (Task 1a), so a step size that is stable along the
          large-scale features is far too small for the others. Accordingly, the condition number of
          $\mathbf X^\top\mathbf X$ drops from $1.8\cdot10^{9}$ on the raw data to $2.9\cdot10^{3}$ on the
          z-scores. With the settings of Task 2 (500 epochs, $\varepsilon_0=0.01$, bold driver), GD reaches
          risk $0.2145$ and training accuracy $92.2\%$ on z-scores, but only risk $2.38$ and accuracy
          $49.7\%$ on the raw data.
    \item \emph{No bias.} Our model has no bias term (Task 2b), so the argument above does not cover shifts:
          z-scoring changes which decision functions the model can represent. Centering the features at $0$
          partly compensates for the missing bias.
    \item \emph{Regularization and interpretation.} The $L_2$ penalty of Task 4 is not invariant to
          rescaling, since features on a small scale need large weights and are therefore penalized more.
          Weights can also only be compared (Task 3c) when all features are on the same scale.
  \end{itemize}
\end{answer}
\section*{Task 2e)}
\begin{answer}[19.5cm]
  \small
  \answerfig{task2_gd_sgd.png}{Settings: $\mathbf w_0\sim\mathcal N(\mathbf 0,\mathbf I)$
  (seed 0), $\varepsilon_0=0.01$, 500 epochs, bold driver. $\hat R^*\approx0.2047$ is estimated from a
  GD run with 30\,000 epochs.}

  \resizebox{\linewidth}{!}{%
  \begin{tabular}{lcccccc}
    \toprule
     & $\hat R$ after 500 ep. & epochs to $1.1\,\hat R^*$ & time/epoch & train acc. & test acc. & $\lVert\mathbf w\rVert_2$ \\
    \midrule
    GD  & 0.2145 & 144 & 0.09\,ms & 92.2\% & 91.9\% & 8.9 \\
    SGD & 0.2164 & 8   & 4.6\,ms  & 92.1\% & 91.7\% & 6.8 \\
    \bottomrule
  \end{tabular}}

  \textbf{Per epoch (a).} SGD makes far more progress early on. Its risk is $0.32$ after one epoch, and
  it comes within $10\%$ of $\hat R^*$ after 8 epochs. This is expected, since an SGD epoch performs
  $N=3065$ updates while a GD epoch performs one. GD barely improves over roughly the first 50 epochs
  and needs 144 epochs to reach the same level.

  \textbf{Step sizes (c).} GD's slow start is a step-size problem: with $\varepsilon_0=0.01$, a single
  full-gradient step per epoch moves $\mathbf w$ very little. The bold driver increases the step size
  by $5\%$ per epoch until it reaches about $60$ in epoch 179. After that it oscillates between roughly
  $8$ and $35$ (23 halvings), and GD overtakes SGD. For SGD, the bold driver is counterproductive.
  Gradient noise from single examples often raises the risk from one epoch to the next, which triggers
  36 halvings. The step size decays to about $10^{-3}$, and SGD stalls at a risk of $0.216$. Three
  further seeds give the same final risk to within $\pm0.00002$, so the random example order has a
  negligible effect here.

  \textbf{Per running time (b).} In wall-clock time, GD is clearly faster. A GD epoch consists of
  vectorized matrix--vector products, whereas an SGD epoch is a Python loop over 3065 examples and takes
  about 50 times as long. GD reaches a lower risk after $0.04$\,s than SGD does after $2.3$\,s, i.e.\
  after all 500 epochs. Part of this gap comes from the implementation rather than the algorithms; a
  vectorized or compiled SGD loop would be considerably faster per epoch.

  \textbf{Convergence.} Neither method has converged after 500 epochs, and both remain about $0.01$
  above $\hat R^*$. In the 30\,000-epoch reference run, the risk continues to decrease slowly while
  $\lVert\mathbf w\rVert_2$ grows to about $50$. For prediction, the remaining gap matters little. Both
  models reach about $92\%$ training and test accuracy, and the final GD and SGD weight vectors have a
  correlation of $0.96$.

  \textbf{Conclusion.} On this small data set, full-batch GD is the better choice. Its epochs are cheap,
  and the bold driver lets it exploit large step sizes. SGD becomes worthwhile when the full gradient
  is expensive to compute, i.e.\ for very large $N$. To converge, however, it needs a decreasing
  step-size schedule rather than a heuristic that reacts to risk increases caused by gradient noise.
\end{answer}
\section*{Task 3b) and 3c)}
\begin{answer}[19.5cm]
Both models were evaluated after 500 epochs on the same test set of 1,536 emails, with 3,065 emails used for training and the same seeded initialization. GD achieved 91.86\% accuracy and a mean test log-loss of 0.2792, while SGD achieved 91.73\% accuracy and 0.2827. The confusion matrices, with rows representing true labels and columns representing predicted labels, are GD $\begin{bmatrix}887&54\\71&524\end{bmatrix}$ and SGD $\begin{bmatrix}886&55\\72&523\end{bmatrix}$. For spam classification, GD achieved a precision of 0.9066 and a recall of 0.8807, compared with 0.9048 and 0.8790 for SGD. The precision--recall curves are nearly identical, and precision remains above the test-set spam prevalence baseline (0.387) across much of the recall range.

The largest-magnitude coefficients correspond to \texttt{word\_freq\_3d}, \texttt{capital\_run\_length\_longest}, and \texttt{char\_freq\_\$}, all with positive weights, while \texttt{word\_freq\_hp} has a negative weight. Among the smallest coefficients are \texttt{word\_freq\_650}, \texttt{word\_freq\_make}, and \texttt{word\_freq\_report}. Positive coefficients increase the spam logit, whereas negative coefficients decrease it. These weights reflect associations learned by the models rather than causal relationships; moreover, correlated features may capture overlapping predictive information.

\begin{center}
\includegraphics[width=0.57\linewidth]{results/task3/confusion_matrices.png}

\includegraphics[width=0.45\linewidth]{results/task3/precision_recall.png}\hfill
\includegraphics[width=0.45\linewidth]{results/task3/feature_weights.png}
\end{center}
\end{answer}

\section*{Task 4b)}
\begin{answer}[19.5cm]
We trained L2-regularized logistic regression with gradient descent for 500 epochs at each
$\lambda\in\{0,0.001,0.01,0.1,1,10,100,1000\}$, always starting from the same initial vector. The
objective minimized during training is $\mathrm{ER}(w)+\tfrac{\lambda}{2}\|w\|_2^2$, but the reported
training risk is the unregularized mean log-loss $\mathrm{ER}(w)$.

Log-loss and accuracy favour different values of $\lambda$. Training risk increases monotonically
from 0.21445 at $\lambda=0$ to 0.69267 at $\lambda=1000$, as expected, since the penalty pulls $w$
away from the minimizer of the empirical risk. Test log-loss rises over the same range from 0.27915 to
0.69267. Test accuracy, in contrast, first improves slightly: from 91.86\% at $\lambda=0$ it reaches a
maximum of 92.06\% at $\lambda=0.01$, then falls to 91.15\% at $\lambda=0.1$ and 89.32\% at $\lambda=1$,
and levels off at 88.74\% for $\lambda=10$ and 88.61\% for $\lambda=100$ and $1000$.

The gain at $\lambda=0.01$ is only 0.20 percentage points, and the test log-loss of this model
(0.29038) is higher than that of the unregularized one. The two metrics disagree because they measure
different things. Accuracy only checks on which side of the threshold each prediction falls, whereas
log-loss scores the predicted probabilities and penalizes confident mistakes heavily. A model can
therefore classify slightly more emails correctly while assigning less accurate probabilities overall.
Since we only stored aggregate metrics, we cannot identify which individual predictions changed. The
choice of $\lambda$ thus depends on the evaluation metric: log-loss favours $\lambda=0$, accuracy
$\lambda=0.01$.

For large $\lambda$, both training and test risk approach $\log 2\approx0.693$, the loss of predicting
0.5 for every example. The predicted probabilities then carry little information, and the model
clearly underfits.
\begin{center}
\includegraphics[width=0.94\linewidth]{results/task4/regularization_metrics.png}
\end{center}
\end{answer}
\section*{Task 4c)}
\begin{answer}[19.5cm]
The training objective is $\mathrm{ER}(w)+\tfrac{\lambda}{2}\|w\|_2^2$, so increasing $\lambda$ penalizes
large coefficients more strongly. Because the model has no bias term, every component of $w$ is
penalized. The weight plot follows the eight coefficients with the largest absolute values at
$\lambda=0$. Their magnitudes generally shrink as $\lambda$ grows, and so does the norm of the full weight
vector: $\|w\|_2$ falls from 8.9448 at $\lambda=0$ to 0.0628 at $\lambda=10$ and 0.000692 at
$\lambda=1000$. The three largest coefficients at $\lambda=0$, for \texttt{word\_freq\_3d},
\texttt{capital\_run\_length\_longest} and \texttt{char\_freq\_\$} (6.9897, 2.3443 and 2.2043), are
0.0000, 0.0001 and 0.0002 at $\lambda=1000$ when rounded to four decimals. \texttt{word\_freq\_hp} keeps
its negative sign but likewise approaches zero.

The shrinkage is not uniform across features. At $\lambda=0.01$, \texttt{word\_freq\_3d} has already
dropped to 0.2622, while \texttt{char\_freq\_\$} retains 1.1009, so the ranking of the coefficients
changes. \texttt{word\_freq\_857} even changes sign between $\lambda=0.01$ and $\lambda=0.1$.
Regularization therefore does more than rescale the unregularized solution: it also changes the
relative weight the model assigns to individual features.

As $w\to\mathbf 0$, the logits $Xw$ tend to zero as well. Without a bias term, the predicted
probabilities $\sigma(Xw)$ then approach 0.5 for every example, and training and test log-loss rise
towards $\log 2$, the loss of a constant prediction of 0.5. Accuracy is affected far less, because the
predicted class depends only on which side of the decision threshold $x^\top w$ falls, not on its
magnitude. At $\lambda=1000$ the model still reaches 88.61\% accuracy, even though its predicted
probabilities are so close to 0.5 that they carry almost no information about confidence.

\begin{center}
\includegraphics[width=0.91\linewidth]{results/task4/weight_composition.png}
\end{center}
\end{answer}
\newpage
\section*{Ehrenwörtliche Erklärung}

Ich versichere, dass ich die beiliegende Bachelor-, Master-, Seminar-,
oder Projektarbeit ohne Hilfe Dritter und ohne Benutzung anderer als der
angegebenen Quellen und in der untenstehenden Tabelle angegebenen
Hilfsmittel angefertigt und die den benutzten Quellen wörtlich oder
inhaltlich entnommenen Stellen als solche kenntlich gemacht habe. Ich
kenne den Inhalt dieser Arbeit und bin in der Lage, ihn auf Nachfrage zu
erklären. Diese Arbeit hat in gleicher oder ähnlicher Form noch keiner
Prüfungsbehörde vorgelegen. Ich bin mir bewusst, dass eine falsche
Erklärung rechtliche Folgen haben wird.

% Declare below which AI tools you used in the process of writing your work,
% including text, image, code, and data generation.
%
% The AI declaration must include at least one an entry for every graded subtask
% (i.e., where points are specified). The entry should state whether or not you
% used AI to solve parts of the task. Using AI for things such as language
% improvements or API search does not count as "Solved with AI", but things such
% as code gen, writing drafts, or solving math tasks does.
%
% We may grade tasks that you (partially or fully) solved with AI more
% stringently. So make sure that whenever you use AI, the results are correct.
% Our recommendation is to not use AI to solve the tasks, as this undermines the
% learning goals of this course and, ultimately, your own education.
\begin{center}
  \textbf{Declaration of Used AI Tools} \\[.3em]

  \begin{tabularx}{\textwidth}{lclXc}
    \toprule
    Task & Solved by AI? & Tool & Purpose & Useful? \\
    \midrule
    1a & No & ChatGPT & Code review and correctness checking & \\
    1b & No & ChatGPT & Code review and correctness checking & \\
    1c & No & ChatGPT & Code review and correctness checking & \\
    2a & & & & \\
    2b & & & & \\
    3a & & & & \\
    3b & & & & \\
    4 & & & & \\
    5a & Yes & ChatGPT & Code refinement, debugging, testing & yes \\
    5b & Yes & ChatGPT & Code refinement, debugging, testing, and academic language editing & yes \\
    6a & Yes & ChatGPT & Mathematical derivation checking, code review, and academic language editing \\
    6b & Yes & ChatGPT & Mathematical derivation checking, code review, and academic language editing & yes \\
    6c & Yes & ChatGPT & Mathematical derivation checking, code review, and academic language editing & yes \\
    \bottomrule
  \end{tabularx}
\end{center}

\vspace{2cm}
\noindent Unterschrift\\
\noindent Mannheim, den 25. Oktober 2026\hfill

%%% Local Variables:
%%% mode: LaTeX
%%% TeX-master: "a01-nb"
%%% End:


\end{document}
